\begin{lemma}
For integers $4\leq s\leq k$, there is a loopless $T_s$-free digraph $D$
with
\[
 N=2^{2s-3}-2^{s-1}-2^{s-2}+1
\]
vertices and
\[
 \operatorname{fwi}_k(D)\leq
 \sum_{t=1}^{s-1}\binom kt,2^{(s-1)(t+k)-\binom{t+1}{2}}.
\]
The sum intentionally counts ambient vector tuples, so it may overcount
projective points and nonzero representatives.
\end{lemma}
\begin{proof}
Put $p=s-1$, identify the points of $G(s-2,2)$ with the nonzero vectors of
$\mathbb F_2^p$, and let $D$ have vertices $(x,y)$ with
$\langle x,y\rangle=1$, with an arc $(x,y)\to(x',y')$ exactly when
$\langle x,y'\rangle=0$.  The polarity count gives
$|V(D)|=(2^p-1)(2^{p-1}-1)=N$, and the triangular orthogonality obstruction
makes $D$ loopless and $T_s$-free.

Write a forward-independent $k$-tuple as
$(a_1,b_1,\ldots,a_k,b_k)$.  It is represented by an ambient, possibly zero,
vector tuple satisfying
$\langle a_j,b_i\rangle=1$ whenever $j\leq i$; counting these ambient tuples
is an upper bound and is the deliberate overcount in the statement.  Let
$r_i=\dim\operatorname{span}(a_1,\ldots,a_i)$ and $r_0=0$.  A realizable rank
sequence has $r_1=1$, each increment is either $0$ or $1$, and its final value
is some $t\leq p$.  There are at most $\binom kt$ such sequences.

Fix one with final rank $t$.  Whenever the rank increases from $i-1$ to $i$,
there are at most $2^p$ choices for $a_i$ and at most $2^{p-i}$ choices for
$b_i$.  Whenever it does not increase, there are at most $2^i$ choices for
$a_i$ and at most $2^{p-i}$ choices for $b_i$.  Multiplying in the order of
the rank increases gives at most
\[
 \left(\prod_{i=1}^t2^{2p-i}\right)2^{p(k-t)}
 =2^{pt+pk-\binom{t+1}{2}}
\]
tuples with this rank sequence.  Summing over the at most $\binom kt$ valid
sequences and then over $1\leq t\leq p=s-1$ gives the displayed bound.
\end{proof}
